\chapter*[Conclusão]{Conclusão}\addcontentsline{toc}{chapter}{Conclusão}\label{conclusao}

Estimulada por artistas-programadores(as), a interiorização de um processo criativo (bricolagem) pode ter conduzido à formação de programas de pesquisas interdisciplinares, entre computação, cognição e práticas musicais tonais. Isto é, teorias das Artes Músicais, Audiovisuais, Corporais, Têxteis, Filosofia, e se a criatividade permitir, da Engenharia Crítica, são apropriadas para elaboração de conceitos. Uma classe particular de artistas-programadores(as) \cite[p.~16]{McLean2011}, o(a) improvisador(a) de códigos (\emph{live coder}), admite conhecimentos da engenharia de computação, revisada sob o prisma de uma ou mais teorias como extensão para improvisações ou ou análise de registros. O improvisador de códigos, em geral, pratica o \emph{live coding} segundo uma cultura de publicações, sejam elas divulgadas na \emph{internet}\disponivelem{http://www.toplap.org}, deselvolvida a partir de \cite{ward_live_2004}, em revistas de programas de investigação científica, como  \emph{Computer Music Journal}, \emph{Leonardo}, em listas de email de \emph{softwares}. Por fim, é estimulado codificar programas, não como um produto de venda, mas para os mais diversos fins. \citeonline{ward_live_2004} definem a improvisação de códigos como \traducao{atividade da escrita integral (ou partes) de um programa enquanto ele é executado}{Live coding is the activity of writing (parts of ) a program while it runs}. \citeonline{blackwell_programming_2005} enfatizam a definição do ponto de vista da linguagem de programação como instrumento musical. \citeonline{mclean_hacking_2006} relata o \emph{live coding} como ferramenta para um \emph{Disk Jockey codificado}.  \citeonline{sorensen_keith_2009} pontuam a improvisação de códigos como \traducao{uma prática de performance para o qual linguagens de computador definem o meio primário de expressão artística.}{Live coding is a performance pratice for which computer languages define  the primary means of expression.}. Para \citeonline{sorensen_impromptu_2010}, \emph{live coding} (ou \emph{livecoding}) envolve a premissa de uma programação-partitura audiovisual reativa: 

\traduzcitacao{
Livecoding é uma prática de arte computacional que envolve criação em tempo-real de programas de audiovisual generativo para performances multimídias interativas. Comumente as ações dos programadores são expostas para uma audiência por projeção do ambiente de edição. Performances de livecoding geralmente envolvem mais de um participante, e são geralmente iniciadas a partir de uma folha conceitual em branco   
\cite[p.~823]{sorensen_impromptu_2010}}{Livecoding [10, 50] is a computational arts practice that involves the real-time creation of generative audiovisual software for interactive multimedia performance. Commonly the programmers’ actions are exposed to the audience by projection of the editing environment. Livecoding performances often involve more than one participant, and are often commenced from a conceptual blank slate}

\citeonline{magnusson_algorithms_2011,collins_origins_2014} sintetizam o \emph{live coding} como improvisação audiovisual. \citeonline{sorensen_programming_2014} define como \traducao{programar sistemas de tempo-real em tempo real}{programming real-time systems in real-time}. Em um discussão entitulada ``\emph{Wtf is livecoding}''\disponivelem{http://lurk.org/groups/livecode/messages/topic/ofAxZpxsKFpDRLnoA48Bh} diz que o \traducao{\emph{Live coding} celebra a efemeridade da própria definição}{Live coding celebrates the ephemerality of definition itself} \ver{fig:live_coding_def}. Embora semelhantes, as definições mudam de detalhes de acordo com o contexto. Por exemplo, \citeonline{ward_live_2004} enfatizam que o código pode ser (re)composto de partes menores. \citeonline{mclean_hacking_2006} pontua o ambiente de codificação onde o código é (re)programado, e que tipo de música é feita. \citeonline{sorensen_programming_2014} destaca que modificar alguma coisa (bricolagem) é próprio da técnica, de forma que é possível extender essa bricolagem para ideologias, enfraquecendo a própria definição de \emph{live coding}. Nick \citeonline{collins_origins_2014} situa essa questão da seguinte formafundamentos objeto:

  \begin{figure}[h]
    \centering
    \includegraphics[scale=0.7]{imagens/live_coding_def.png}
    \caption{Definição de \emph{live coding}: ``Insira a definição aqui''. \textbf{Fonte}: \citeonline{collins_origins_2014}.}
    \label{fig:live_coding_def}
  \end{figure}


Se por um lado a definição agrega definições, o que dificulta a tarefa inicial de descrever os fundamentos do objeto de pesquisa, por outro ilustra a improvisação de códigos como um \emph{Universo de conceitos}. Neste trabalho consideramos que definições ou performances de improvisação de códigos estão contidas em diferentes \emph{Espaços Conceituais} \cite{wiggins_framework_2006,mclean_music_2006}. Artistas-programadores (\emph{live coders}) transitam entre os Espaços Conceituais para criação de Sistemas Criativos (códigos, programas). Estes Sistemas Criativos são representados em diferentes Linguagens de Programação. Regras práticas conduzem o processo de escrita e exposição desta linguagem; mas não restringem o resultado (no caso da pesquisa, musical). Mas algumas categorizações musicais se destacam.  Neste sentido, selecionamos um exemplo simbólico, \emph{A Study in Keith} de Andrew \citeonline{sorensen_keith_2009}\footnote{Disponível em \url{https://vimeo.com/2433947}.}. Representa um caso particular que foge dos exemplos citados anteriormente, mas envolve a manutenção de uma tradição musical tonal através de um interessante esforço de \emph{replicação do estilo}. No \autoref{cap:introducao} selecionamos alguns exemplos afim de ilustrar nossa percepção (conhecimento-psicológico) do imaginário daquilo que \citeonline{McLean2011} chama de artistas-programadores. No \autoref{sec:protohistoria} buscamos levantar um conjunto de conhecimentos históricos. No \autoref{cap:metodologia}, discutimos um modelo de formalização conceitos, observados pelo prisma de Alex \citeonline{mclean_music_2006}, contraposto por \citeonline{thornton_quantitative_2007}, e rediscutido por \citeonline{Forth2010} \citeonline{McLean2011}. No \autoref{cap:estudos_de_caso}, organizamos conceitos de um algoritmo inicial de uma improvisação de códigos, \emph{A Study in Keith} de 2009, segundo este um Quadro Conceitual de Sistemas Criativos. Fnalizamos este trabalho com o apêndice \autoref{app:A}, onde descrevemos um processo de organização de tais qualidades da improvisação de códigos. 